% Einheit 2 von 6 – Nullstellen (bank/funktionsklassen-und-eigenschaften/e2.jsonl, zusammenbau v0.1)
%% TODO Zeitmarke Einheit 2: Marken-Zeile nicht lesbar
%% TODO Zweigzeile Teil 1: was hier gelernt wird (Satz in Schülersprache aus dem Einheitstitel) – Einheit 2
\einheitenkopf[e2]{Einheit 2 von 6 · Nullstellen}
\zweigzeile{Abitur GK}
\setcounter{aufgabe}{11}

%% TODO Titel: Ich-kann-Satz fehlt in der Bank (Platzhalter: Kettenname) – Nr. 12
%% TODO Anweisung: ein Satz über der ersten Teilaufgabe fehlt in der Bank – Nr. 12
\begin{aufgabe}{Nullstellen}
%% TODO \rechenplatz{n}: Schrittzahl des Grundfalls fehlt in der Bank (nur bei fester Schreibform, 2.3 b)
\begin{teile}
\teil $f(x) = (x + 4)(x - 7)$ – welches Verfahren? \\ \kreuz{Faktoren null setzen} \\ \kreuz{ausklammern} \\ \kreuz{substituieren} \\ \kreuz{Polynomdivision} \\ \kreuz{Lösungsformel}
\teil $f(x) = (x - 2)(x + 5)$ – Nullstellen? x1 = \leerfeld , x2 = \leerfeld
\end{teile}
\begin{gleichungsraster}[2]
\gl{\text{$f(x) = x^3 - 9x$ – Nullstellen?}} &
\gl{\text{$f(x) = x^2 - 6x + 5$ – Nullstellen?}} \\
\gl{\text{$f(x) = x^4 - 13x^2 + 36$ – Nullstellen?}} &
\gl{\text{$f(x) = x^3 - 2x^2 - 5x + 6$ – bestätige die Nullstelle $1$ und bestimme alle Nullstellen.}} \\
\end{gleichungsraster}
\begin{teile}
\teil $f(x) = (x^2 - 4) \cdot e^x$ – Nullstellen?
\teil $f(x) = x^3 - x^2 - 6x$ – alle Schnittpunkte mit den Achsen?
\teil Ein Tunnelquerschnitt wird durch $f(x) = -0{,}5x^2 + 4{,}5$ beschrieben (1 Längeneinheit entspricht 1 m, Boden auf der x-Achse). Breite des Tunnels am Boden?
\end{teile}
\begin{gleichungsraster}[2]
\gl{\text{$f(x) = 0{,}2x^4 - 2x^2 + 1{,}8$ – berechne die Koordinaten aller Schnittpunkte des Graphen mit den Koordinatenachsen (FHR 2026).}} & \\
\end{gleichungsraster}
\end{aufgabe}

%% TODO Titel: Ich-kann-Satz fehlt in der Bank (Platzhalter: Kettenname) – Nr. 13
%% TODO Anweisung: ein Satz über der ersten Teilaufgabe fehlt in der Bank – Nr. 13
\begin{aufgabe}{Nullstellen und Werte: y-Achsenschnittpunkt angeben und Anzahl der Nullstellen aus Grenzverhalten und Tiefpunkt ohne Rechnung begründen}
\begin{teile}
\teil $f(x) = x^4 - 4x + 5$ hat genau einen Tiefpunkt, $T(1 | 2)$. Schnittpunkt mit der y-Achse? Wie viele Nullstellen hat $f$? Begründe ohne Rechnung.
\end{teile}
\end{aufgabe}

%% TODO Titel: Ich-kann-Satz fehlt in der Bank (Platzhalter: Kettenname) – Nr. 14
%% TODO Anweisung: ein Satz über der ersten Teilaufgabe fehlt in der Bank – Nr. 14
\begin{aufgabe}{Nullstellen – Fehler finden, Begründen, Anwendung}
\begin{teile}
\teil Ich finde den Fehler: \rechnung{x^3 - 16x &= 0 &&\mid :x \\ x^2 &= 16 \\ x &= \pm 4}
\teil Begründe: $(x - 5)(x + 1) = 0$ gilt genau für $x = 5$ und $x = -1$.
\teil Beim Kugelstoßen gilt $h(x) = -0{,}05x^2 + 0{,}8x + 1{,}8$ ($h$ Höhe in m, $x$ waagerechte Entfernung vom Abstoß in m). Wie weit fliegt die Kugel?
\end{teile}
\end{aufgabe}
